\documentclass[11pt,a4paper]{article}
\usepackage[margin=1in]{geometry}
\usepackage{graphicx}
\usepackage{amsmath}
\usepackage{caption}
\usepackage{subcaption}
\usepackage{tikz}
\usepackage{hyperref}

\title{Deep Learning Methods for Tiger Video Generation}
\author{Your Name}
\date{}

\begin{document}
\maketitle

\section{Data}

\subsection{Tiger Image Dataset}
To generate realistic tiger videos, high-resolution images of tigers are collected from wildlife footage and open-source datasets such as ImageNet-Vid and YouTube-VIS. Images include frontal, profile, and tilted views with diverse lighting conditions.  

Preprocessing steps:
\begin{itemize}
    \item Resize all images to $128\times128$ pixels,
    \item Normalize pixel values to fit model distribution,
    \item Data augmentation: flipping, color jitter, and rotation to increase variability.
\end{itemize}

\subsection{Motion Data (Video Clips)}
Videos of tigers exhibiting different movements (walking, running, turning) are extracted. Optical flow or keypoint-based motion representations are computed to provide temporal guidance.  

\subsection{Audio (Optional)}
If considering sound-driven motion (e.g., roar lip-sync), audio is processed into Mel-Spectrograms and temporally aligned with video frames.

---

\section{Network Architecture}

\subsection{GAN Baseline: MoCoGAN-style}
The baseline framework uses a GAN to generate video clips from a random noise vector.

\textbf{Generator:}
\begin{itemize}
    \item Inputs: content latent vector $z_c$ (appearance) and motion latent vector $z_m(t)$ (temporal dynamics),
    \item 3D convolutional layers or ConvLSTM to generate video sequences.
\end{itemize}

\textbf{Discriminator:}
\begin{itemize}
    \item 2D CNN for per-frame realism,
    \item 3D CNN for temporal consistency.
\end{itemize}

% \begin{figure}[h]
%     \centering
%     % Placeholder for GAN flowchart
%     \includegraphics[width=0.7\textwidth]{gan_flowchart.png}
%     \caption{GAN baseline framework for video generation.}
% \end{figure}

\subsection{Diffusion-based Video Generator (Main)}
Diffusion models generate temporally coherent tiger videos via iterative denoising in latent space.

\textbf{Architecture:}
\begin{itemize}
    \item Input: random noise video clip + motion condition,
    \item 3D U-Net with temporal attention for denoising,
    \item Optional latent VAE encoder/decoder for efficiency.
\end{itemize}

\textbf{Losses:}
\begin{itemize}
    \item Score matching loss for denoising,
    \item Perceptual loss to preserve texture and details,
    \item Optional temporal consistency loss.
\end{itemize}

% \begin{figure}[h]
%     \centering
%     % Placeholder for Diffusion flowchart
%     \includegraphics[width=0.7\textwidth]{diffusion_flowchart.png}
%     \caption{Diffusion-based video generation pipeline.}
% \end{figure}

\subsection{Joint Motion Conditioning}
For more controllable outputs, motion embeddings extracted from optical flow or pose keypoints can be injected into intermediate layers via adaptive normalization, ensuring realistic movement while preserving identity.

---

\section{Resources \& Infrastructure}

\subsection{Pre-trained Models \& Tools}
\begin{itemize}
    \item GAN: MoCoGAN pre-trained weights (optional),
    \item Diffusion: pre-trained Video Diffusion Model or VAE backbone,
    \item Motion extractor: RAFT or keypoint detector,
    \item Audio: Librosa for Mel-Spectrogram extraction.
\end{itemize}

\subsection{Hardware \& Data Pipeline}
High-performance GPUs (e.g., NVIDIA A100 or RTX 4090) are recommended. Efficient data loaders (PyTorch DataLoader) support large-scale video batches.

\subsection{Evaluation Metrics}
\begin{itemize}
    \item Video quality: FVD, PSNR, SSIM, LPIPS,
    \item Identity preservation: cosine similarity via pre-trained recognition model (ArcFace),
    \item Temporal consistency: Average Pose Distance (APD) or frame-to-frame motion correlation.
\end{itemize}

---

\section{Design Philosophy \& Thoughts}

\begin{itemize}
    \item \textbf{Modular design:} Separate GAN baseline and diffusion modules for clear comparison.
    \item \textbf{Temporal consistency:} Diffusion with temporal attention ensures smooth motion.
    \item \textbf{Motion-content disentanglement:} Motion conditions are injected at intermediate layers, preserving spatial fidelity.
    \item \textbf{Stage-wise training:} Pre-train modules individually before joint fine-tuning.
    \item \textbf{Future extensions:} Text- or audio-conditioned tiger generation for controllable video synthesis.
\end{itemize}

\end{document}